% Summary for the rerun-rust profile: Software Engineer (Rust), Rerun (Remote Europe); see profiles/rerun-rust.tex for the not-claimed list.
\IfLanguageName{english}{
    % Cut for the one-page fit; the thesis and agent clauses, kept: next to the Vulkan and OpenGL renderers I have maintained since 2020; my B.Sc.\ thesis was on image quality in real-time rendering. At Fraunhofer IPA I build camera pipelines and GPU services, maintain repositories in Rust, C++ and Python, and started the institute's self-hosted coding agent.
    % Work experience leads (owner, 2026-11-15); the summary before, kept: I write real-time graphics and GPU software in Rust and C++. In my own time I build a wgpu renderer in Rust with an egui overlay that runs natively and in the browser, next to the Vulkan and OpenGL renderers I have maintained since 2020. At Fraunhofer IPA I build camera pipelines and GPU services in Rust, C++ and Python. Coding agents are part of how I write and review code every day, and fixes I need in other people's tools go upstream, as they did for LLVM and sccache.
    \par{At Fraunhofer IPA I build real-time camera pipelines, GPU services and the infrastructure under them in Rust, C++ and Python: activity detection for a Horizon Europe project, the institute's self-hosted coding agent on vLLM, and the CI/CD and build systems the group works from. I maintain several of the institute's repositories, and fixes I need in other people's tools go upstream, as they did for LLVM and sccache. I come from computer graphics and build a wgpu renderer in Rust on the side.}
}{
    % Für die eine Seite gekürzt; die Fassung davor, erhalten: neben den Vulkan- und OpenGL-Renderern, die ich seit 2020 pflege; meine Bachelorarbeit behandelte Bildqualität im Echtzeit-Rendering. Am Fraunhofer IPA baue ich Kamera-Pipelines und GPU-Dienste, pflege Repositories in Rust, C++ und Python und habe den selbst gehosteten Coding-Agent des Instituts gestartet.
    % Berufserfahrung zuerst (Inhaber, 2026-11-15); die Fassung davor, erhalten: Ich schreibe Echtzeitgrafik- und GPU-Software in Rust und C++. In meiner Freizeit entwickle ich einen wgpu-Renderer in Rust mit egui-Overlay, der nativ und im Browser läuft, neben den Vulkan- und OpenGL-Renderern, die ich seit 2020 pflege. Am Fraunhofer IPA baue ich Kamera-Pipelines und GPU-Dienste in Rust, C++ und Python. Coding-Agents gehören täglich dazu, wenn ich Code schreibe und prüfe, und Korrekturen, die ich in fremden Werkzeugen brauche, gehen upstream, so bei LLVM und sccache.
    \par{Am Fraunhofer IPA baue ich Echtzeit-Kamera-Pipelines, GPU-Dienste und die Infrastruktur darunter in Rust, C++ und Python: Aktivitätserkennung für ein Horizon-Europe-Projekt, den selbst gehosteten Coding-Agenten des Instituts auf vLLM sowie die CI/CD- und Build-Systeme, mit denen die Gruppe arbeitet. Ich pflege mehrere Repositories des Instituts, und Korrekturen, die ich in fremden Werkzeugen brauche, gehen upstream, so bei LLVM und sccache. Ich komme aus der Computergrafik und entwickle nebenbei einen wgpu-Renderer in Rust.}
}
